\chapter{Functions, Domain and Range}

Up to here every chapter has manipulated \emph{expressions}. A
\de{Funktion} is the next step: an expression together with a decision
about which numbers may be substituted into it, and a record of which
numbers come out. Those two sets, the \de{Definitionsbereich} (domain) and
the \de{Wertebereich} (range), are not decoration. A limit is taken at a
point of the domain or at its boundary; a derivative exists only where the
function does; an optimisation problem is solved over a domain; and a
demand function that returns a negative price has left the region where it
means anything.

Finding a maximal domain is pure algebra of the kind the previous chapters
drilled: a denominator must not vanish, an even root needs a non-negative
radicand, a logarithm needs a positive argument. Each condition is an
equation or an inequality, and the domain is what survives all of them at
once. Finding a range is slightly different work: either the function is
built from pieces whose ranges are known, or the equation $y=f(x)$ is
solved for $x$ and one asks which $y$ admit a solution.

\section{Core Rules}

\begin{rulebox}{Function, domain, range}
A function $f\colon D\to\R$ assigns to every $x\in D$ exactly one value
$f(x)\in\R$. Two sets belong to the description:
\begin{itemize}[leftmargin=1.6em,itemsep=0.3ex]
  \item the \textbf{domain} $D(f)$ \de{(Definitionsbereich)}: the $x$ that
        may be substituted;
  \item the \textbf{range} $W(f)=\{f(x): x\in D(f)\}$
        \de{(Wertebereich, Bildmenge)}: the values actually attained.
\end{itemize}
The \textbf{graph} of $f$ is the set $\{(x,f(x)) : x\in D(f)\}$ in the
plane. A curve is the graph of a function exactly when every vertical line
meets it at most once (\textbf{vertical line test}). Reading off a graph:
$D(f)$ is its shadow on the $x$-axis, $W(f)$ its shadow on the $y$-axis.
\end{rulebox}

\begin{gallery}
\begin{tikzpicture}
  \begin{axis}[fnaxis, xmin=-0.6, xmax=5.4, ymin=-0.6, ymax=3.4,
      xtick={1,3,5}, ytick={1,2,3}, title={$f(x)=\sqrt{x-1}$}]
    \addplot[domain=1:5]{sqrt(x-1)};
  \end{axis}
\end{tikzpicture}\hfill
\begin{tikzpicture}
  \begin{axis}[fnaxis, xmin=-3.4, xmax=3.4, ymin=-0.6, ymax=4.4,
      xtick={-2,2}, ytick={2,4}, title={$f(x)=x^{2}$}]
    \addplot[domain=-2.1:2.1]{x^2};
  \end{axis}
\end{tikzpicture}\hfill
\begin{tikzpicture}
  \begin{axis}[fnaxis, xmin=-3.4, xmax=3.4, ymin=-1.2, ymax=4.4,
      xtick={-2,2}, ytick={1,3}, title={$f(x)=\mathrm{e}^{x}$}]
    \addplot[domain=-3:1.45]{exp(x)};
  \end{axis}
\end{tikzpicture}
\gcap{$D=[1,\infty)$, $W=[0,\infty)$ \quad\textbullet\quad
$D=\R$, $W=[0,\infty)$ \quad\textbullet\quad $D=\R$, $W=(0,\infty)$.}
\end{gallery}

\begin{rulebox}{The maximal domain: three obstructions}
Unless a domain is prescribed, $D(f)$ is the \textbf{maximal domain}: every
real $x$ for which the formula produces a real number. Only three things
can go wrong, and each is an (in)equality to solve:
\begin{itemize}[leftmargin=1.6em,itemsep=0.3ex]
  \item \textbf{Denominator:} $\dfrac{\;\cdot\;}{g(x)}$ requires
        $g(x)\neq 0$. Solve $g(x)=0$ and delete those points.
  \item \textbf{Even root:} $\sqrt[2k]{g(x)}$ requires $g(x)\geq 0$.
        Odd roots ($\sqrt[3]{\;}$) are defined for all reals.
  \item \textbf{Logarithm:} $\ln g(x)$ requires $g(x)>0$ (strictly).
\end{itemize}
A power $g(x)^{r}$ with non-integer $r$ counts as a root: it requires
$g(x)\geq 0$ for $r>0$ and $g(x)>0$ for $r<0$. If several obstructions
occur, \textbf{intersect} the resulting sets. A root in a denominator
combines two conditions into one: $\dfrac{1}{\sqrt{g(x)}}$ needs
$g(x)>0$.
\end{rulebox}

\begin{rulebox}{Domains prescribed by context}
In economics the domain is usually narrower than the maximal domain and
comes from the model, not from the formula:
quantities and inputs satisfy $q\geq 0$, $K,L\geq 0$; a price must satisfy
$p\geq 0$; a share or probability lives in $[0,1]$; an average cost
$AC(q)=C(q)/q$ needs $q>0$. State the economic domain explicitly, because
the formula will happily return values outside it: the linear demand
$p=120-3q$ is defined for every real $q$, but only $q\in[0,40]$ keeps both
$q\geq0$ and $p\geq 0$.
\end{rulebox}

\begin{rulebox}{Finding the range}
Four routes, in increasing order of generality:
\begin{enumerate}[leftmargin=1.7em,itemsep=0.3ex]
  \item \textbf{Standard building blocks.}
        $x^{2}\colon[0,\infty)$, \;$\sqrt{x}\colon[0,\infty)$,
        \;$|x|\colon[0,\infty)$, \;$\mathrm{e}^{x}\colon(0,\infty)$,
        \;$\ln x\colon\R$, \;$1/x\colon\R\setminus\{0\}$,
        \;$x^{3}\colon\R$, \;$\sin x,\cos x\colon[-1,1]$.
        Then track what the outer operations do: adding $c$ shifts the
        range by $c$, multiplying by $a>0$ scales it, multiplying by $a<0$
        scales \emph{and} flips it.
  \item \textbf{Complete the square} for a quadratic:
        $a(x-h)^{2}+k$ has range $[k,\infty)$ if $a>0$ and
        $(-\infty,k]$ if $a<0$.
  \item \textbf{Monotonicity.} If $f$ is strictly increasing on $[a,b]$,
        then $W=[f(a),f(b)]$; open ends stay open, and an end at
        $\pm\infty$ is replaced by the limiting value.
  \item \textbf{Solve $y=f(x)$ for $x$} (the general method). The range is
        the set of $y$ for which the solution exists and lies in $D(f)$.
        For $f(x)=\dfrac{x}{x+1}$: \;$y(x+1)=x\Rightarrow x=\dfrac{y}{1-y}$,
        solvable for every $y\neq 1$, hence $W=\R\setminus\{1\}$.
\end{enumerate}
\end{rulebox}

\begin{rulebox}{Composite functions}
$(g\circ f)(x)=g\bigl(f(x)\bigr)$ is evaluated \textbf{inside out}. Its
domain carries two conditions at once:
\[
  D(g\circ f)=\{\,x\in D(f)\;:\;f(x)\in D(g)\,\}.
\]
So $x$ must be admissible for the inner function, and its \emph{output}
must be admissible for the outer one. For $g(u)=\ln u$ and $f(x)=x^{2}-1$:
$x^{2}-1>0$, hence $D=(-\infty,-1)\cup(1,\infty)$. Order matters:
$g\circ f\neq f\circ g$ in general, and the two domains usually differ.
\end{rulebox}

\begin{rulebox}{Symmetry, monotonicity, injectivity}
\begin{itemize}[leftmargin=1.6em,itemsep=0.3ex]
  \item $f$ is \textbf{even} \de{(gerade)} if $f(-x)=f(x)$ for all $x\in D$
        (graph symmetric about the $y$-axis), \textbf{odd}
        \de{(ungerade)} if $f(-x)=-f(x)$ (symmetric about the origin).
        Only even powers $\Rightarrow$ even; only odd powers
        $\Rightarrow$ odd; a mix is neither.
  \item $f$ is \textbf{strictly increasing} if $x_{1}<x_{2}$ implies
        $f(x_{1})<f(x_{2})$, strictly decreasing if the inequality
        reverses.
  \item $f$ is \textbf{injective} \de{(eineindeutig)} if $f(x_{1})=f(x_{2})$
        forces $x_{1}=x_{2}$: no horizontal line meets the graph twice.
        Strictly monotone $\Rightarrow$ injective $\Rightarrow$ invertible
        on its range (Chapter~3).
\end{itemize}
\end{rulebox}

\section{Worked Example}

\begin{worked}{Find the maximal domain of $f(x)=\dfrac{\ln(x+3)}{\sqrt{4-x^{2}}}$, and the range of $g(x)=\dfrac{x+3}{x-2}$}
\wstep{Step 1. List the obstructions, one per structure in the formula.}
Three of them appear here, which is why this example was chosen:
\begin{itemize}[leftmargin=1.6em,itemsep=0.2ex]
  \item the logarithm needs $x+3>0$;
  \item the even root needs $4-x^{2}\ge0$;
  \item the denominator needs $\sqrt{4-x^{2}}\neq0$.
\end{itemize}
\wtrap{The last two combine into the \emph{strict} inequality $4-x^{2}>0$.
A square root sitting in a denominator is not allowed to be zero, so the
usual ``$\ge$'' for an even root tightens to ``$>$''. Handling the root and
the denominator as separate conditions and then forgetting to intersect them
is the standard way to hand in $[-2,2]$ instead of the right answer.}

\wstep{Step 2. Solve each condition separately.}
\[
  x+3>0 \iff x>-3 ;
  \qquad
  4-x^{2}>0 \iff x^{2}<4 \iff -2<x<2 .
\]

\wstep{Step 3. Intersect --- never unite.}
The function is defined where \emph{all} conditions hold at once:
\[
  D_{f}=(-3,\infty)\cap(-2,2)=\boxed{(-2,2)} .
\]
\wtrap{``And'' means intersection. Each obstruction removes points; nothing
in the formula ever gives points back. If your answer is larger than one of
the individual conditions, you took a union by mistake.}

\wstep{Step 4. Check the boundary behaviour, which confirms the strictness.}
As $x\to2^{-}$ the denominator $\to0^{+}$ while $\ln5>0$, so $f\to+\infty$.
The endpoint really must be excluded. At $x=-2$ the same thing happens.

\wstep{Step 5. The range of $g$: solve $y=g(x)$ for $x$.}
The range is the set of $y$ for which the equation $y=g(x)$ has a solution.
So treat $y$ as given and solve:
\[
  y=\frac{x+3}{x-2}
  \;\Longrightarrow\;
  y(x-2)=x+3
  \;\Longrightarrow\;
  yx-x=2y+3
  \;\Longrightarrow\;
  x(y-1)=2y+3 .
\]

\wstep{Step 6. Read off which $y$ fail.}
\[
  x=\frac{2y+3}{y-1},
\]
which exists for every $y$ except $y=1$; and at $y=1$ the previous line reads
$0=5$, impossible. Hence
\[
  \boxed{\;\text{range}(g)=\R\setminus\{1\}\;}
\]

\wstep{Step 7. Confirm it against the graph.}
$g(x)=\frac{x+3}{x-2}=1+\frac{5}{x-2}$, so the graph is the hyperbola
$5/x$ shifted, with horizontal asymptote $y=1$. The value $1$ is approached
but never attained --- exactly what the algebra said. Two independent
routes agreeing is worth the extra line.
\wtrap{Do not assume the range of a rational function is all of $\R$, and do
not assume it is the same as the domain. The excluded input ($x=2$) and the
excluded output ($y=1$) are different numbers, read off different parts of
the formula.}
\end{worked}

\section{Practice Problems}

State the maximal domain $D(f)\subseteq\R$ in problems 1--18, using
interval notation.

\begin{problems}
  \item $f(x)=\dfrac{x+3}{x-2}$
  \item $f(x)=\sqrt{x-5}$
  \item $f(x)=\sqrt{9-x^{2}}$
  \item $f(x)=\ln(x-4)$
  \item $f(x)=\ln(x^{2}-4)$
  \item $f(x)=\dfrac{\sqrt{x}}{x-1}$
  \item $f(x)=\dfrac{1}{\sqrt{x-3}}$
  \item $f(x)=\sqrt{x^{2}-x-6}$
  \item $f(x)=\dfrac{\ln x}{x-2}$
  \item $f(x)=\mathrm{e}^{1/x}$
  \item $f(x)=\sqrt{x+2}+\sqrt{4-x}$
  \item $f(x)=\dfrac{x^{2}-1}{x^{2}-3x+2}$
  \item $f(x)=\ln(\ln x)$
  \item $f(x)=x^{3/4}$
  \item $f(x)=\dfrac{1}{1-\mathrm{e}^{x}}$
  \item $f(x)=\sqrt{\ln x}$
  \item $f(x)=\sqrt[3]{x-7}$
  \item $f(x)=\dfrac{\sqrt{4-x}}{x^{2}-x}$
\end{problems}

\begin{problems}[resume]
  \item A firm has cost $C(q)=50+2q$. State the economic domain of $C$
        and of the average cost $AC(q)=C(q)/q$.
  \item For the demand schedule $p(q)=120-3q$, state the largest interval
        of quantities on which both $q\geq0$ and $p\geq0$ hold.
  \item State the domain of the Cobb--Douglas product
        $Y(K,L)=K^{0.3}L^{0.7}$ as a subset of the $(K,L)$-plane, and say
        why the exponents force the restriction.
  \item Utility is $U(x)=\ln x$. State the economic domain and explain what
        happens to $U$ as $x\to 0^{+}$.
\end{problems}

Determine the range $W(f)$ in problems 23--36. The domain is the maximal
one unless stated otherwise.

\begin{problems}[resume]
  \item $f(x)=x^{2}+4$
  \item $f(x)=x^{2}-6x+11$
  \item $f(x)=-2x^{2}+8x-3$
  \item $f(x)=\mathrm{e}^{x}-3$
  \item $f(x)=5-\mathrm{e}^{x}$
  \item $f(x)=\ln(x+2)$
  \item $f(x)=\dfrac{1}{x-2}+3$
  \item $f(x)=\sqrt{x-1}$
  \item $f(x)=\sqrt{16-x^{2}}$
  \item $f(x)=|x-3|$
  \item $f(x)=\dfrac{2x+1}{x-3}$
  \item $f(x)=\dfrac{x^{2}}{x^{2}+1}$
  \item $f(x)=x^{2}$ on $D=[-1,3]$
  \item $f(x)=2x-5$ on $D=(1,4]$
\end{problems}

\begin{problems}[resume]
  \item With $f(u)=\sqrt{u}$ and $g(x)=x-4$: give $(f\circ g)(x)$ and
        $(g\circ f)(x)$ with both domains.
  \item With $f(u)=1/u$ and $g(x)=x+2$: give $(f\circ g)(x)$ and its
        domain.
  \item With $f(u)=\ln u$ and $g(x)=9-x^{2}$: give $(f\circ g)(x)$ and its
        domain.
  \item With $f(u)=\sqrt{u}$ and $g(x)=1-x^{2}$: give $(f\circ g)(x)$, its
        domain and its range.
  \item Classify as even, odd, or neither:
        (a) $x^{4}-3x^{2}$, \;(b) $x^{3}-x$, \;(c) $x^{2}+x$,
        \;(d) $\dfrac{1}{x}$.
  \item Which of $x^{3}$, $x^{2}$, $\mathrm{e}^{x}$, $|x|$ are injective on
        $\R$?
  \item Give domain and range of
        $f(x)=\begin{cases} x^{2}, & x<1,\\ 2x-1, & x\geq 1.\end{cases}$
  \item Does $x^{2}+y^{2}=25$ define $y$ as a function of $x$? Justify with
        the vertical line test, and give the two functions obtained by
        solving for $y$, with their domains.
  \item A graph consists of the segment from $(-2,1)$ to $(0,5)$ together
        with the horizontal ray from $(0,5)$ to the right. Read off
        $D$ and $W$.
\end{problems}

\section{Solutions}

\begin{solutions}
  \item Denominator $x-2\neq0$: \;$D=\R\setminus\{2\}
        =(-\infty,2)\cup(2,\infty)$.
  \item $x-5\geq0$: \;$D=[5,\infty)$.
  \item $9-x^{2}\geq0\iff x^{2}\leq9\iff|x|\leq3$: \;$D=[-3,3]$.
  \item $x-4>0$: \;$D=(4,\infty)$.
  \item $x^{2}-4>0\iff|x|>2$: \;$D=(-\infty,-2)\cup(2,\infty)$.
  \item Two conditions: $x\geq0$ for the root and $x\neq1$ for the
        denominator: \;$D=[0,1)\cup(1,\infty)$.
  \item The root sits in the denominator, so $x-3>0$ (not merely
        $\geq0$): \;$D=(3,\infty)$.
  \item $x^{2}-x-6=(x-3)(x+2)\geq0$; the parabola opens upwards, so the
        product is non-negative outside the roots:
        \;$D=(-\infty,-2]\cup[3,\infty)$.
  \item $\ln x$ needs $x>0$, the denominator needs $x\neq2$:
        \;$D=(0,2)\cup(2,\infty)$.
  \item Only $x\neq0$; the exponential itself is unrestricted:
        \;$D=\R\setminus\{0\}$.
  \item $x+2\geq0$ and $4-x\geq0$, i.e. $x\geq-2$ \emph{and} $x\leq4$:
        \;$D=[-2,4]$.
  \item $x^{2}-3x+2=(x-1)(x-2)$: \;$D=\R\setminus\{1,2\}$. (The numerator
        also vanishes at $x=1$, but a cancelling factor does not repair the
        domain: $x=1$ stays excluded.)
  \item Inner: $x>0$; outer: $\ln x>0\iff x>1$. The stronger condition
        wins: \;$D=(1,\infty)$.
  \item Exponent $3/4$ is not an integer, so the base must satisfy
        $x\geq0$: \;$D=[0,\infty)$.
  \item $1-\mathrm{e}^{x}=0\iff \mathrm{e}^{x}=1\iff x=0$:
        \;$D=\R\setminus\{0\}$.
  \item $\ln x$ needs $x>0$, and $\ln x\geq0$ needs $x\geq1$; the
        intersection is \;$D=[1,\infty)$.
  \item An odd root accepts every real radicand: \;$D=\R$.
  \item Root: $4-x\geq0\iff x\leq4$. Denominator: $x^{2}-x=x(x-1)\neq0$,
        so $x\neq0$ and $x\neq1$. Together:
        \;$D=(-\infty,0)\cup(0,1)\cup(1,4]$.
  \item $C$: production cannot be negative, so $D(C)=[0,\infty)$; at $q=0$
        the value $C(0)=50$ is the fixed cost. $AC$: dividing by $q$
        excludes the origin, so $D(AC)=(0,\infty)$.
  \item $p\geq0\iff 120-3q\geq0\iff q\leq40$; with $q\geq0$ this gives
        $q\in[0,40]$.
  \item $D=\{(K,L): K\geq0,\;L\geq0\}$, the non-negative quadrant. The
        exponents $0.3$ and $0.7$ are not integers, so $K^{0.3}$ and
        $L^{0.7}$ require non-negative bases; negative input quantities
        would be meaningless anyway.
  \item $D=(0,\infty)$: consumption must be strictly positive, since
        $\ln 0$ is undefined. As $x\to0^{+}$, $U(x)\to-\infty$: the
        logarithmic utility punishes zero consumption infinitely hard,
        which is exactly why interior solutions appear in Cobb--Douglas
        problems.
  \item $x^{2}\geq0$, so $x^{2}+4\geq4$: \;$W=[4,\infty)$.
  \item $x^{2}-6x+11=(x-3)^{2}+2$: \;$W=[2,\infty)$.
  \item $-2x^{2}+8x-3=-2(x^{2}-4x)-3=-2(x-2)^{2}+5$; the parabola opens
        downwards: \;$W=(-\infty,5]$.
  \item $\mathrm{e}^{x}\in(0,\infty)$, shifted down by $3$:
        \;$W=(-3,\infty)$.
  \item $\mathrm{e}^{x}\in(0,\infty)$, negated to $(-\infty,0)$, shifted up
        by $5$: \;$W=(-\infty,5)$.
  \item $\ln$ attains every real value, and $x+2$ runs through all of
        $(0,\infty)$ on $D=(-2,\infty)$: \;$W=\R$.
  \item $1/(x-2)$ attains every value except $0$; adding $3$ shifts this:
        \;$W=\R\setminus\{3\}$.
  \item A square root returns exactly the non-negative numbers, and
        $x-1$ runs through $[0,\infty)$: \;$W=[0,\infty)$.
  \item $D=[-4,4]$, on which $16-x^{2}$ runs through $[0,16]$; the root of
        that is \;$W=[0,4]$.
  \item An absolute value is non-negative and attains every non-negative
        value (at $x=3\pm w$): \;$W=[0,\infty)$.
  \item Solve $y=\dfrac{2x+1}{x-3}$: \;$y(x-3)=2x+1\Rightarrow
        x(y-2)=3y+1\Rightarrow x=\dfrac{3y+1}{y-2}$, which exists exactly
        for $y\neq2$: \;$W=\R\setminus\{2\}$. ($y=2$ is the horizontal
        asymptote.)
  \item Solve $y=\dfrac{x^{2}}{x^{2}+1}$: \;$y(x^{2}+1)=x^{2}\Rightarrow
        x^{2}=\dfrac{y}{1-y}$. A real $x$ exists iff
        $\dfrac{y}{1-y}\geq0$, i.e. $0\leq y<1$: \;$W=[0,1)$.
  \item On $[-1,3]$ the square is not monotone: it falls to $0$ at $x=0$
        and rises to $9$ at $x=3$. Minimum $0$, maximum $9$:
        \;$W=[0,9]$.
  \item $f$ is strictly increasing, so the endpoints map to the endpoints:
        $f(1)=-3$ (excluded, as $1\notin D$) and $f(4)=3$ (included):
        \;$W=(-3,3]$.
  \item $(f\circ g)(x)=\sqrt{x-4}$ with $D=[4,\infty)$;
        $(g\circ f)(x)=\sqrt{x}-4$ with $D=[0,\infty)$. Same ingredients,
        different functions and different domains.
  \item $(f\circ g)(x)=\dfrac{1}{x+2}$, requiring $x+2\neq0$:
        \;$D=\R\setminus\{-2\}$.
  \item $(f\circ g)(x)=\ln(9-x^{2})$, requiring $9-x^{2}>0\iff|x|<3$:
        \;$D=(-3,3)$.
  \item $(f\circ g)(x)=\sqrt{1-x^{2}}$, requiring $1-x^{2}\geq0$:
        \;$D=[-1,1]$. On it $1-x^{2}$ runs through $[0,1]$, so
        $W=[0,1]$ (the upper unit semicircle).
  \item (a) even (only even powers); (b) odd (only odd powers);
        (c) neither: $f(-x)=x^{2}-x$ equals neither $f(x)$ nor $-f(x)$;
        (d) odd: $1/(-x)=-1/x$.
  \item $x^{3}$ and $\mathrm{e}^{x}$ are injective (both strictly
        increasing). $x^{2}$ and $|x|$ are not: $f(-1)=f(1)$ in both cases.
  \item $D=\R$. First piece: $x$ runs through $(-\infty,1)$, so $x^{2}$
        attains every value in $[0,\infty)$. Second piece: $2x-1$ on
        $[1,\infty)$ gives $[1,\infty)$. Union: \;$W=[0,\infty)$.
  \item No. The vertical line $x=0$ meets the circle at $(0,5)$ and
        $(0,-5)$, so the assignment is not unique. Solving gives two
        functions, $y=\sqrt{25-x^{2}}$ (upper semicircle) and
        $y=-\sqrt{25-x^{2}}$ (lower), each with domain $[-5,5]$ and ranges
        $[0,5]$ and $[-5,0]$.
  \item $D=[-2,\infty)$: the segment covers $x\in[-2,0]$, the ray covers
        $x\in[0,\infty)$. On the segment $y$ climbs from $1$ to $5$, on the
        ray $y$ stays at $5$: \;$W=[1,5]$.
\end{solutions}

\section{Additional Problems}

\begin{problems}
  \item $f(x)=\dfrac{3x}{2x+7}$: maximal domain.
  \item $f(x)=\sqrt{2x+8}$: maximal domain and range.
  \item $f(x)=\ln(5-2x)$: maximal domain.
  \item $f(x)=\dfrac{1}{\ln x}$: maximal domain.
  \item $f(x)=\dfrac{x-1}{\sqrt{x^{2}-4}}$: maximal domain.
  \item $f(x)=\sqrt{\dfrac{x-1}{x+2}}$: maximal domain.
  \item $f(x)=\ln\!\bigl(x^{2}+1\bigr)$: maximal domain and range.
  \item $f(x)=\mathrm{e}^{-x^{2}}$: maximal domain and range.
  \item $f(x)=\dfrac{1}{x^{2}-5x+6}$: maximal domain.
  \item $f(x)=\sqrt{4-x}+\dfrac{1}{x+1}$: maximal domain.
  \item $f(x)=(x-3)^{-1/2}$: maximal domain.
  \item $f(x)=\dfrac{|x|}{x}$: maximal domain and range.
  \item $f(x)=3-\dfrac{2}{x+1}$: range.
  \item $f(x)=\dfrac{x-1}{x+1}$: range.
  \item $f(x)=4-\sqrt{x+9}$: range.
  \item $f(x)=x^{2}-4x+7$ on $D=[0,3]$: range.
  \item $f(x)=\mathrm{e}^{2x}+1$ on $D=[0,\ln 2]$: range.
  \item $f(x)=\ln x$ on $D=[1,\mathrm{e}^{3}]$: range.
  \item With $f(u)=u^{2}$ and $g(x)=\sqrt{x-3}$: give $(f\circ g)(x)$ and
        its domain, and explain why the simplified formula alone would give
        the wrong answer.
  \item With $f(u)=\dfrac{1}{u-1}$ and $g(x)=\mathrm{e}^{x}$: give
        $(f\circ g)(x)$ and its domain.
  \item Revenue is $R(q)=q\,(200-4q)$. State the economic domain on which
        both $q\ge 0$ and the price stay non-negative, and the range of $R$
        on it.
  \item A production function is $Q(L)=10\sqrt{L}$ with $L\in[0,100]$.
        Give the range, and the average product $Q(L)/L$ with its domain.
  \item Classify as even, odd, or neither: (a) $\mathrm{e}^{-x^{2}}$,
        \;(b) $x\,|x|$, \;(c) $\ln x$, \;(d) $\dfrac{x^{3}}{x^{2}+1}$.
  \item Show that $f(x)=\dfrac{2x-1}{x+4}$ is injective on its domain by
        solving $f(x_{1})=f(x_{2})$.
  \item Give domain and range of
        $f(x)=\begin{cases} -x, & x<0,\\ \sqrt{x}, & 0\le x\le 4,\\
        2, & x>4.\end{cases}$
\end{problems}

\section{Additional Solutions}

\begin{solutions}
  \item $2x+7\neq0\iff x\neq-\tfrac{7}{2}$:
        \;$D=\R\setminus\{-\tfrac{7}{2}\}$.
  \item $2x+8\geq0\iff x\geq-4$: \;$D=[-4,\infty)$, and since $2x+8$ then
        runs through $[0,\infty)$, \;$W=[0,\infty)$.
  \item $5-2x>0\iff x<\tfrac{5}{2}$: \;$D=(-\infty,\tfrac{5}{2})$.
  \item $x>0$ for the logarithm and $\ln x\neq0\iff x\neq1$:
        \;$D=(0,1)\cup(1,\infty)$.
  \item The root is in the denominator, so $x^{2}-4>0\iff|x|>2$:
        \;$D=(-\infty,-2)\cup(2,\infty)$.
  \item Need $\dfrac{x-1}{x+2}\geq0$ with $x\neq-2$. The quotient is
        non-negative outside the interval between the two critical points
        $-2$ and $1$: \;$D=(-\infty,-2)\cup[1,\infty)$. (At $x=1$ the
        quotient is $0$, which is allowed; $x=-2$ is not.)
  \item $x^{2}+1\geq1>0$ always, so $D=\R$; the argument runs through
        $[1,\infty)$, hence \;$W=[0,\infty)$.
  \item $D=\R$. The exponent $-x^{2}$ runs through $(-\infty,0]$, and
        $\mathrm{e}^{u}$ maps that to $(0,1]$: \;$W=(0,1]$ (the bell curve,
        maximum $1$ at $x=0$).
  \item $x^{2}-5x+6=(x-2)(x-3)$: \;$D=\R\setminus\{2,3\}$.
  \item $4-x\geq0\iff x\leq4$, and $x\neq-1$:
        \;$D=(-\infty,-1)\cup(-1,4]$.
  \item A negative non-integer exponent means a root in a denominator:
        $x-3>0$, \;$D=(3,\infty)$.
  \item $D=\R\setminus\{0\}$. The quotient equals $1$ for $x>0$ and $-1$
        for $x<0$: \;$W=\{-1,1\}$, a range consisting of two points.
  \item $2/(x+1)$ attains every value except $0$; negating and adding $3$
        keeps exactly one value out of reach: \;$W=\R\setminus\{3\}$.
  \item $y=\dfrac{x-1}{x+1}\Rightarrow y(x+1)=x-1\Rightarrow
        x(y-1)=-1-y\Rightarrow x=\dfrac{-(1+y)}{y-1}=\dfrac{1+y}{1-y}$,
        which exists for $y\neq1$: \;$W=\R\setminus\{1\}$.
  \item $D=[-9,\infty)$; $\sqrt{x+9}$ runs through $[0,\infty)$, and
        $4-\sqrt{x+9}$ therefore through \;$W=(-\infty,4]$.
  \item $x^{2}-4x+7=(x-2)^{2}+3$. On $[0,3]$ the vertex $x=2$ lies inside,
        giving the minimum $3$; the endpoints give $f(0)=7$ and $f(3)=4$,
        so the maximum is $7$: \;$W=[3,7]$.
  \item Strictly increasing: $f(0)=2$ and
        $f(\ln 2)=\mathrm{e}^{2\ln 2}+1=4+1=5$: \;$W=[2,5]$.
  \item Strictly increasing: $\ln 1=0$ and $\ln \mathrm{e}^{3}=3$:
        \;$W=[0,3]$.
  \item $(f\circ g)(x)=\bigl(\sqrt{x-3}\bigr)^{2}=x-3$, but only where the
        inner function is defined: \;$D=[3,\infty)$. The simplified formula
        $x-3$ suggests $D=\R$; the domain is inherited from the inner
        function and is not restored by cancellation.
  \item $(f\circ g)(x)=\dfrac{1}{\mathrm{e}^{x}-1}$, requiring
        $\mathrm{e}^{x}\neq1\iff x\neq0$: \;$D=\R\setminus\{0\}$.
  \item Price $200-4q\ge0\iff q\le50$, so the domain is $[0,50]$.
        $R(q)=200q-4q^{2}=-4(q-25)^{2}+2500$, a downward parabola with
        vertex at $q=25$; on $[0,50]$ it runs from $0$ up to $2500$ and
        back to $0$: \;$W=[0,2500]$.
  \item $\sqrt{L}$ runs through $[0,10]$ on $[0,100]$, so
        $W=[0,100]$. The average product is
        $Q(L)/L=10/\sqrt{L}$, which needs $L>0$: domain $(0,100]$.
  \item (a) even: $(-x)^{2}=x^{2}$; (b) odd:
        $(-x)|-x|=-x|x|$; (c) neither, and in fact the test cannot even be
        applied: the domain $(0,\infty)$ is not symmetric about $0$;
        (d) odd: numerator odd, denominator even.
  \item $\dfrac{2x_{1}-1}{x_{1}+4}=\dfrac{2x_{2}-1}{x_{2}+4}$
        $\Rightarrow(2x_{1}-1)(x_{2}+4)=(2x_{2}-1)(x_{1}+4)$
        $\Rightarrow 2x_{1}x_{2}+8x_{1}-x_{2}-4
        =2x_{1}x_{2}+8x_{2}-x_{1}-4$
        $\Rightarrow 9x_{1}=9x_{2}\Rightarrow x_{1}=x_{2}$. Injective on
        $\R\setminus\{-4\}$.
  \item $D=\R$. First piece: $-x$ on $(-\infty,0)$ gives $(0,\infty)$.
        Second: $\sqrt{x}$ on $[0,4]$ gives $[0,2]$. Third: the constant
        $2$. The union is \;$W=[0,\infty)$.
\end{solutions}
